\section{The problem}
\label{sec:problem}

A 4-billion-parameter model in f16 holds 8.04~GB of weights. Quantized to two
bits it is 1.79~GB on disk. Whether the GPU also \emph{reads} two bits per
weight is a separate question, and for the Leech lattice the answer has been
no. The format we published earlier stores 2.000 bits of code per weight and
its kernel reads 4.804 bits per weight from VRAM (Fig.~\ref{fig:gap}).
Batch-1 decoding is bound by weight traffic, so the second number is the one
that sets the cost.

The gap comes from the size of the codebook. \citet{llvq2026} quantize 24
weights at a time on $\Lambda_{24}$, with a 47-bit index into a ball of
$1.1\times10^{14}$ lattice points and one gain bit: 48 bits per block, 2.000
bits per weight. Fast 2-bit kernels decode such a label with a lookup table.
QTIP keeps a 16-bit trellis state and reads 2~KiB~\citep{qtip2024}; QuIP\#
reads an $E_8$ codebook of $2^{16}$ entries~\citep{quipsharp2024}; AQLM and
VPTQ learn codebooks sized for a table~\citep{aqlm2024,vptq2024}. A table with
$1.1\times10^{14}$ entries is not an option. We measured the two alternatives
in earlier work~\citep{llvq1preprint}. Decoding the combinatorial index
arithmetically inside the kernel was compute-bound in every variant we built,
so we shipped the other one: unfolding the index at load time into a wide
shift-and-mask stream. That is where 4.804 comes from.

\begin{figure}[t]
\centering
\includegraphics[width=\linewidth]{fig_gap.pdf}
\Description{Four horizontal dumbbells, one per format. Each joins a hollow
marker at the bits of code the format carries per weight to a filled marker at
the bits its kernel reads per weight in VRAM. Planes14 runs from 2.000 to
4.804, Tetra from 2.000 to 2.148, QTIP has no gap, AWQ runs from 4.000 to
4.179.}
\caption{What a format carries against what its kernel reads, per weight:
hollow marker to filled marker, circles for our arms and squares for the
deployed kernels. All three 2-bit formats carry 2.000 bits of code.
\Planes{} reads $2.40\times$ that, because its codebook has to be unfolded
first. \Tetra{} reads $1.07\times$ it; the excess is the exactly-kept tail and
the row scales, not the code. FP16 is omitted: 16.000 at both ends. Data:
echelle-formats.csv.}
\label{fig:gap}
\end{figure}

This paper keeps the lattice and changes the codebook. Reordering the 24
coordinates into three disjoint octads of the Golay code turns the code into a
trellis with 64 states per cut. A lattice point then has a short address: an
8-bit state and three 11-bit rows of one 16~KiB table that all three sections
share. The word is still 48 bits, the kernel reads it as stored, and the VRAM
rate falls to 2.148 bits per weight. Sectionalising $\mathcal{G}_{24}$ is a
classical construction~\citep{forney1988,vardy1991}; what we propose is its use
as the address space of a quantizer codebook, together with the truncation rule
it forces and the shaping that rule gives up.

\S\ref{sec:tetra} builds the codebook and states what it costs.
\S\ref{sec:kernel} is the kernel and one measurement about where its time
goes. \S\ref{sec:results} measures a single Qwen3-4B object: bytes, time,
memory, perplexity, accuracy. \S\ref{sec:limitations} says what one model on
one card does not support.
